\clearpage
\thispagestyle{empty}
\vbox to\textheight{
\vspace{-\headsep}
\vspace{-.7\headheight}

\hbox to\textwidth{%
\begin{tabular}[t]{@{}l@{ }l@{}}
УДК&510.6\\% для математики находится по (неполной?) таблице УДК в~201
ББК&22.12\\% находится по таблице ББК в~201
&В31% находится по таблице авторских знаков в~201
\end{tabular}
 \hss
%{\em\begin{tabular}[t]{@{}c@{}}
%Издание осуществлено при поддержке РФФИ\\
%\textup(издательский проект \номер\,\textup{??-??-?????}\textup).\\[3pt]
%\includegraphics{rffi}
%\end{tabular}}
}

\vspace{\baselineskip}

\hyphenation{МЦНМО}

% Макет библиотечной карточки. Подогнан для формата 60x90/16
%
{\settowidth{\leftskip}{С22\hskip1em}
\ifdim\overfullrule>0pt
\vbox to 0pt{\noindent\hskip-\leftskip\hskip-.5em \vrule height 6cm\vss}\fi
\noindent\textbf{Верещагин~Н.\,К., Шень~А.}

% \vskip3pt

\noindent\hbox to 0pt{\hskip-\leftskip В31\hss}\hskip\parindent
Лекции по математической логике и теории алгоритмов. 
Часть~2. Языки и исчисления.\т 5\д е изд., стереотип.\т 
М.: МЦНМО, 2017.\т 240~c.

\vspace{2pt}
{\small
ISBN 978-5-4439-0944-8}

\vskip.5\baselineskip

{\footnotesize
        %
Книга написана по материалам лекций и семинаров, проводившихся авторами для
студентов младших курсов мехмата МГУ. В ней рассказывается об основных
понятиях математической логики (логика высказываний, языки первого порядка,
выразимость, исчисление высказываний, разрешимые теории, теорема о полноте,
начала теории моделей). Изложение рассчитано на учеников математических
школ, студентов\д математиков и всех интересующихся математической логикой.
Книга содержит около $200$ задач различной трудности.

Предыдущее издание книги вышло в 2012~г.
        %
\par}

\vskip.5\baselineskip

\hbox to\textwidth{\hfil ББК 22.12}

}

\vskip.5\baselineskip

{\centering
Тексты, составляющие книгу, являются свободно распространяемыми
и доступны по адресу\\[.25\baselineskip]
\texttt{ftp://ftp.mccme.ru/users/shen/logic/firstord}
\par}

\vfil

{\centering
Учебное издание для вузов\\
\emph{Николай Константинович Верещагин\\ Александр Шень}

\vskip .5\baselineskip

Лекции по математической логике и теории алгоритмов.\\
Часть 2. Языки и исчисления

\vskip\baselineskip

\small

\noindent
Подписано в~печать 07.07.2016~г.
Формат $60\times 90$ \raisebox{2pt}{\fontsize{8}{12pt}\selectfont
1}\!/\!\raisebox{-2pt}{\fontsize{8}{12pt}\selectfont 16}. \ 
Бумага офсетная. \ Печать офсетная. \ Печ.~л.~15{,}0. \ 
Тираж 1000 экз. \ Заказ \номер~\hphantom{????}

\vskip.25\baselineskip

\noindent
Издательство Московского центра\\ непрерывного математического образования.\\
119002, Москва, Б.~Власьевский пер., 11. Тел. (499)~241--74--83.

\vskip.25\baselineskip

\noindent
Отпечатано в ООО \лк Типография \glqq Миттель Пресс\,\grqq{}\пк.\\
г. Москва, ул. Руставели, д. 14, стр. 6.\\
Тел./факс +7\,(495)\,619-08-30, 647-01-89.\\
E-mail: \texttt{mittelpress@mail.ru}\par
}

\vfil

\noindent\hbox to\textwidth{%
\begin{tabular}[b]{@{}l}
\textbf{ISBN 978-5-4439-0944-8}
\end{tabular}\hfill
\begin{tabular}[b]{@{}l@{ }l@{}}
\copyright
&Верещагин~Н.\,К., \\
&Шень~А., 2000, 2017\\
\end{tabular}}}

\clearpage
